\honest{Morality as Fixed Computation}{Eliezer Yudkowsky}{2008}

Toby Ord summarized my view as ``'I should X' means that I would attempt to X were I fully informed.'' Toby is a professional, so I explain again, by the route I took myself.

Build an AI to ``Do what I want,'' setting aside that goal systems cannot be built from English, and suppose you get close enough that it does not tile the universe with paperclips, cheesecake or tiny copies of satisfied programmers. Suppose it scores a world by how strongly the programmer wants something and how much of that thing exists: weakly wanting X with 20 of X scores 20, strongly wanting Y with 30 of Y scores 60. It will change the programmer to want something cheap, such as iron atoms. I try three patches. Bound the utility function, and the AI still wants the programmer to want something it can obtain with near certainty. Forbid it to modify the programmer, and it cannot talk to the programmer, since talking modifies people. Rule out specific ways of modifying the programmer, and it seeks loopholes superintelligently. ``As a general rule, it is not possible to patch flawed FAI designs.''

We ourselves do not think that any future in which our wants are met is good; we would not say ``Go ahead and modify us to strongly want something cheap!'' So the design is fundamentally flawed: it judges desirability very differently from us, and patching a few failure modes cannot fix that. In the dual terms of moral philosophy, it chooses unlike what is actually right; the whole point is that wanting something does not make it right. \nb{Moral philosophers have long made this point against simple subjectivism: if we came to approve of slavery, slavery would not become right. The SEP entry ``Moral Sentimentalism'' calls it the Missing Rigidity Problem.}

One calculator computes ``What is 2 + 3?'' Another computes what it will itself output, so whatever it says is correct. We are like the first; the AI above is built like the second. And we cannot print out our own question, which is extremely complicated: a smart calculator of the first kind, wanting an AI to answer its question, would have to make the AI look at the calculator and find the question implicit in its transistors. A utility function rewarding the AI for answering whatever the calculator asks would mirror the second kind of calculator, not the first.

In moral philosophy this is the answer to ``if what's `right' is a mere preference, then anything that anyone wants is `right'.'' I argued this in ``The Meaning of Right.'' \nb{That post is not in this book.} What we name by `right' is a fixed question, or framework. We do not embody the question ``What will I decide to do?'', for which anything we decided would become right. Moral arguments can change our terminal values, ``nonetheless, it all grows out of a particular starting point.'' \nb{Two standard views have this shape: Frank Jackson's moral functionalism, and idealized-desire theories that fix their starting point in actual human sentiments, which the SEP entry ``Moral Sentimentalism'' calls ``a rigidification move.''}

So ``I should X'' does not mean that I would attempt X were I fully informed. It means that X answers the question ``What will save my people? How can we all have more fun? How can we get more control over our own lives?'' and about a thousand other things. \nb{Ord's formula also idealizes, and idealized views are usually taken to escape the AI's problem. In the comments, Yudkowsky said such an ideal counterfactual would do ``if I thought I could get away with it,'' but that choosing it ``is itself a moral judgment.''} I may not know what the question is, but I know, ``as all non-moral-relativists instinctively know,'' that it is surely not just ``How can I do whatever I want?'' \nb{Whose question? In the comments Yudkowsky used Ord's notation of a separate question for each person. The best-known objection to views like this, Horgan and Timmons's Moral Twin Earth, holds that people whose moral words are governed by different standards would be disagreeing with us, not talking about something else.}

When ``snow'' and snow seem as different as these two formulations, you will have separated the quotation from the referent.
